\begin{textsample}{POS Dim 2 – pi – Score 31.00 – pi/pi\_em\_23.txt}  \label{ex:f2_pos_094}
4.2 Expectativas dos sujeitos do Ensino Médio De acordo com as orientações nacionais para a (re)construção dos \textbf{currículos}, uma das premissas para êxito no processo seria a obtenção de um diagnóstico da Rede, que permitisse a construção de uma proposta coerente com a perspectiva das escolas, considerando todas as potencialidades e os contextos diversos e adversos.
Além da obtenção de dados concretos, a equipe responsável pela (re)elaboração do \textbf{currículo} deveria conhecer o nível de compreensão da comunidade escolar sobre o “ novo ” previsto pela \textbf{reforma} do Ensino Médio e, ainda, “ escutar ” seus anseios, a fim de se construir uma proposta que estivesse em \textbf{conformidade} com os interesses demonstrados.
Destarte essa orientação, a equipe ProBNCC entendeu que a obtenção desse diagnóstico exigiria um esforço de uma equipe que pudesse atuar para além do processo de escrita.
Então, seguiu -se o caminho inverso : a partir de estudo aprofundado das orientações constantes dos diversos manuais de orientação já citados em momento anterior, iniciou -se o projeto da escrita, com base em informações existentes nos arquivos da SEDUC.
Entretanto, reconheceu -se que os dados não \textbf{atendiam}, pela insuficiência de informações, aos anseios dos redatores e nem, ao menos, davam um direcionamento para a escrita.
Assim, definiu -se inicialmente um recorte de 3 ( três ) escolas por \textbf{Gerência} Regional, como piloto para o envio de um questionário a estudantes e professores, cujas respostas norteassem a construção do \textbf{currículo}, haja vista que todas as questões eram voltadas para as mudanças previstas para o Ensino Médio.
A definição do quantitativo de escolas se deu considerando diferentes modalidades dentre as vinte e uma \textbf{Gerências} Regionais de Educação ( GREs ).
No entanto, no decorrer do processo o questionário foi compartilhado por um número maior de escolas, do qual se obteve retorno da maioria e de algumas não.
De modo que obtivemos respostas de escolas situadas em 24 ( vinte e quatro ) \textbf{municípios}.
A ideia de se aplicar o questionário com o público das escolas selecionadas tinha por finalidade observar se havia coerência das respostas dadas por estudantes e professores com o desenho da arquitetura curricular que a equipe estava propondo para o \textbf{currículo}.
Para grata surpresa, as respostas elencadas estavam em perfeita sintonia com o tom da escrita e o que já se havia desenhado inicialmente e isso proporcionou uma certa segurança na escrita.
Ao se pensar o quantitativo de escolas para as quais seriam enviados os questionários, não se fez o levantamento de quantas respostas seriam obtidas.
Ressalta -se que, ao final do período de escuta, obteve -se o quantitativo de 935 respostas de professores e 4.480 de estudantes.
Esse número ainda é pequeno, considerando a quantidade de escolas de Ensino Médio da Rede.
Mas constitui um corpus significativo para a pesquisa que se pretendia.
Quanto ao \textbf{perfil} dos respondentes, os professores têm idade entre 27 e 40 anos ; residentes, em sua maioria, no mesmo \textbf{município} onde trabalham ( 91,1 \% ) ; mais da metade possui \textbf{curso} superior com \textbf{especialização} ( 64,1 \% ) ; a maior quantidade de respostas foi de professores de Língua Portuguesa ( 17 \% ).
Não foi possível identificar o gênero dos professores porque o questionário não contemplava essa opção.
Já os estudantes, 73,1 \% moram e estudam na zona urbana e 92,1 \% declararam estudar no mesmo \textbf{município} onde residem ; 38,9 \% estão \textbf{cursando} a 1ª série ; 30,9 \% na 2ª e 30,1 \% na 3ª.
Em sua maioria, estudam na modalidade regular de ensino.
Tal qual o questionário dos professores, não se obteve informação acerca do gênero ; também não foi possível identificar a idade dos estudantes.
Para a realização dessa escuta, optou -se pelo questionário desenvolvido pelo MEC ( para as escolas que \textbf{aderiram} ao projeto de flexibilização curricular⁶ ), com alguns ajustes.
Não se utilizou os dados obtidos no questionário aplicado inicialmente com as escolas piloto, por conta da dificuldade de tabulação dos resultados.
Assim, após adequação das questões ao que se pretendia obter como resposta, enviou -se o questionário, por meio da plataforma Google Forms, o que facilitou em demasia o retorno pretendido.
De posse dos dados obtidos, foi feita a sistematização das respostas e estas foram apresentadas às 21 \textbf{Gerências} Regionais de Educação, por meio de um seminário regional, realizado em novembro de 2019, na modalidade presencial ( a equipe ProBNCC foi distribuída em grupos para que todos pudessem participar da atividade ).
A ideia era contextualizar os participantes das mudanças previstas para o Ensino Médio e apresentar o status em que se encontrava a escrita do documento curricular⁷, bem como compartilhar a devolutiva do questionário.
De acordo com as respostas, foi possível perceber que alunos e professores demonstraram conhecimento quanto ao processo de flexibilização curricular, possível a partir dos \textbf{itinerários} \textbf{formativos} e isso deu um suporte à equipe de redatores para a escrita das trilhas de aprendizagem e a proposta de \textbf{eletivas}.
Vale ressaltar que o questionário foi aplicado com estudantes do Ensino Médio – um público que demonstra um pouco mais de \textbf{maturidade} em relação às questões relacionadas à Educação Básica.
Foram apresentadas seis mudanças previstas para o Ensino Médio, dentre as quais destacamos : possibilidade de escolha dos estudantes ; adequação da \textbf{carga} \textbf{horária} ; formação \textbf{profissional} no ensino regular ; existência de uma Base comum.
Destas, a que mais foi marcada como conhecida pelos professores foi a possibilidade de escolha ( 70,1 \% ).
A que foi apontada como menos conhecida foi referente à \textbf{carga} \textbf{horária} ( 39,1 \% dos professores informaram não ter conhecimento prévio sobre o tema ).
O mesmo questionamento foi apresentado aos estudantes.
Eles deviam marcar quais assuntos já eram de conhecimento deles e, tal qual os professores, a possibilidade de fazer escolhas foi a que mais obteve resposta ( 47,1 \% ).
Outro ponto comum foi ao demonstrarem não compreender a proposta de adequação de \textbf{carga} \textbf{horária} ( apenas 16,9 \% demonstraram ter conhecimento ).
Ainda sobre as respostas obtidas, alguns pontos merecem atenção e serão elencados aqui, com as devidas considerações.
Uma das questões que se pretende destacar é : Para você, quais os principais motivos para o aluno \textbf{cursar} o Ensino Médio? ⁸ Como resposta, obteve -se o seguinte resultado : 81,5 \% dos professores responderam que o motivo seria entrar na faculdade. 78,2 \% dos estudantes também deram a mesma resposta, o que leva a crer que um Ensino Médio com foco nos exames nacionais seria mais atrativo.
Entretanto, apesar da resposta apresentada, um paradoxo chamou atenção.
Ao serem questionados sobre quais áreas do conhecimento gostariam de aprofundar ( eles poderiam indicar até 3 ), 42,1 \% dos estudantes escolheram a Educação \textbf{Técnica} e Profissional, como se percebe na descrição abaixo, e isso permite uma inferência de que, talvez, os estudantes do Ensino Médio não tenham \textbf{maturidade} para fazer escolhas. | Área | \textbf{Percentual} | |---|---:| | Matemática e suas Tecnologias | 23,5 \% | | Linguagens e suas Tecnologias | 26,9 \% | | Ciências da Natureza e suas Tecnologias | 32,7 \% | | Ciências Humanas e Sociais Aplicadas | 30,9 \% | | Formação \textbf{Técnica} e Profissional | 42,1 \% | Quadro 1 – Respostas dos estudantes quanto à escolha para aprofundamento Fonte : elaboração própria, a partir dos dados da escuta com os estudantes.
No entanto, a inferência citada no parágrafo anterior é refutada ao se observar o seguinte ponto : indagados sobre se gostariam de fazer algum tipo de Formação \textbf{Técnica} e Profissional ( \textbf{cursos} \textbf{Técnicos} e \textbf{habilitações} \textbf{profissionais} ) durante o Ensino Médio, 78,9 \% responderam sim.
Com base nessas respostas, entende -se que os estudantes esperam, de fato, novidades para o Ensino Médio, mas com uma dinâmica que os permitam construir seus projetos de forma a \textbf{atender} às expectativas em relação a estudos e trabalho, mas sem perder de vista seus sonhos.
Por isso, ao se construir propostas de \textbf{itinerários}, a equipe levou em consideração todas as necessidades que precisavam ser consideradas, embora haja ciência de que muitos ajustes precisam ser feitos para a \textbf{implantação} do \textbf{currículo}.
Importante ressaltar que os resultados obtidos nas escutas estão conectados com o que a equipe de redatores havia pensado como sugestão para a organização curricular, embora o levantamento dos dados tenha ocorrido em etapa posterior ao desenho inicial da proposta.
Destaca -se, no entanto, que uma questão não foi contemplada nessa escuta : a diversidade de contextos coexistentes ( Educação de Jovens e Adultos, Educação Indígena, Educação Quilombola, Educação do Campo, Educação Especial etc. ), por entendimento de que o \textbf{currículo} é geral para todas as modalidades, diferenciando pela forma de \textbf{oferta} e metodologias utilizadas.
Quanto ao aspecto da \textbf{oferta}, os professores foram questionados sobre quais \textbf{itinerários} a escola de atuação deles teria condições de \textbf{ofertar}.
As áreas de Linguagens e Matemática obtiveram, respectivamente, 63,6 \% e 52,2 \% das respostas ; já as áreas de Ciências ( tanto Humanas quanto Natureza ) obtiveram \textbf{percentual} bem menor em relação às duas citadas, ficando as duas com resposta igual ( 39,7 \% ).
Sobre essa percepção dos professores, é válido observar que talvez o desconhecimento da proposta curricular tenha contribuído para obtenção do resultado.
O mesmo se aplica ao Itinerário de Formação \textbf{Técnica} e Profissional : apenas 35,8 \% apontaram possibilidade de \textbf{oferta}.
Outro ponto de atenção é que as respostas dadas por professores talvez sejam diferentes das que dariam os gestores escolares.
A partir das respostas, percebeu -se também que estudantes e professores apresentaram pontos de vista parecidos em relação a alguns tópicos, dentre eles, a sugestão de qual \textbf{profissional} deveria ser o responsável pela condução do componente Projeto de Vida : 53,8 \% dos estudantes responderam “ um \textbf{profissional} especializado ” ; 52 \% dos professores também optaram por essa resposta.
É possível que, sendo este um componente obrigatório somente nas escolas de Tempo Integral, professores e estudantes ainda não tenham uma opinião formada sobre o tema, por não terem vivenciado em sua prática cotidiana.
Há outras questões que não foram analisadas neste tópico, por conta do recorte e da análise que se pretendia.
Mas, de antemão, é possível afirmar que, embora não tenha havido uma conversa prévia com estudantes e professores acerca do Novo ensino Médio, algumas escolas já haviam iniciado esse diálogo, ainda que de forma tímida.
E, quando da realização do Seminário Regional ( mencionado anteriormente ), muitos estudantes e professores manifestaram total interesse e empolgação pela proposta curricular apresentada.
Ainda que haja manifestações contrárias ao estabelecido em Normativos e Resoluções Nacionais.
Quanto à necessidade de obtenção de um diagnóstico, a Rede já iniciou um estudo que permita uma visão ampla de todo o contexto educacional e, por meio dele, seja possível elencar não só as possibilidades de ajustes na arquitetura curricular, como também permita a \textbf{efetivação} de um plano de implementação.
Destaca -se que, no âmbito da Educação Profissional, já existe o desenho das potencialidades da Rede, conforme detalhado neste documento curricular em tópico referente ao Histórico do Ensino Médio no Piauí.
Sobre esse ponto, paralelo ao processo de escrita curricular ( e posterior a ele ), busca -se conhecer a realidade da Rede Estadual quanto ao quantitativo de escolas por Regional ( considerando as diferentes modalidades ) ; professores e alunos por escolas ; infraestrutura ; alimentação ; transporte escolar etc.
Além disso, a Seduc inicia um plano de ação transversal com vistas ao cumprimento de todas as etapas do processo de implementação do \textbf{currículo} para o Novo Ensino Médio.
Destaca -se que, em tópicos seguintes, apresenta -se o desenho de todas as especificidades acerca do DCR do Piauí, iniciando pelo modelo de organização curricular adotado pela Rede, após validação do Comitê de Governança⁹, e que foi apresentado a todas as GREs, em Seminário Regional. ⁶ O questionário foi aplicado nas escolas que \textbf{aderiram} ao Programa de Apoio Novo Ensino Médio, juntos aos alunos, com vistas a \textbf{atender} as \textbf{diretrizes} do MEC de que trata o Documento \textbf{Orientador} do referido programa. ⁷ À época do seminário, toda a parte conceitual do DCR já estava pronta, bem como a definição da arquitetura curricular.
A equipe estava se debruçando sobre a Formação Geral Básica, com a intenção de concluir até o final do referido ano. ⁸ É possível conhecer todo o questionário em arquivo anexo a este documento curricular. ⁹ \textbf{Instituído} pela \textbf{Portaria} \textbf{PORTARIA} SEDUC-PI/GSE Nº 442/2020, de 19 de junho de 2020.

% matched lemmas: aderir, atender, carga, conformidade, currículo, cursar, curso, diretriz, efetivação, eletivo, especialização, formativo, gerência, habilitação, horário, implantação, instituir, itinerário, maturidade, município, oferta, ofertar, orientador, percentual, perfil, portaria, profissional, reforma, técnico
\end{textsample}
